% theme: familiar_substances
\MajorQuestion{身のまわりの物質}
\SetRowSpaceQanda

\Instruction{次の問いに答えなさい。}
\QQ{一定の体積あたりの質量を、その物質の\NamedBlank{density}という。}{密度}
\QQ{ガスバーナーで、炎の大きさを調節するねじを何というか。}{ガス調節ねじ}
\QQ{ものの形や外観に注目したときの呼び名を\NamedBlank{object}という。}{物体}
\QQ{炭素を含み、燃やすと二酸化炭素や水ができる物質を\NamedBlank{organic}という。}{有機物}
\QQ{金、銀、銅、鉄、アルミニウムなどをまとめて\NamedBlank{metal}という。}{金属}

\Instruction{\strut}
\QQ{電子てんびんで粉末状の薬品の質量をはかるとき、薬品をのせる紙を何というか。}{薬包紙}
\QQ{質量54\,g、体積20\,cm$^3$の物質の\RefBlank{density}を求めなさい。}{2.7\,g/cm$^3$}
\QQ{金属を引っ張って細い線状にのばせる性質を何というか。}{延性}
\QQ{ガスバーナーの炎は、空気の量を調節して何色の安定した炎にするか。}{青色}
\QQ{\RefBlank{object}の材料に注目したときの呼び名は何か。}{物質}

\Instruction{\strut}
\QQ{右ききの人が上皿てんびんでものの質量をはかるとき、はかるものは左右どちらの皿にのせるか。}{左の皿}
\QQ{\RefBlank{organic}以外の物質を何というか。}{無機物}
\QQ{メスシリンダーの目盛りを読むとき、目の位置を何と同じ高さにするか。}{液面}
\QQ{ガラス、食塩、プラスチック、木、紙、ゴムなど、\RefBlank{metal}以外の物質を何というか。}{非金属}
\QQ{\RefBlank{density}が7.8\,g/cm$^3$、質量が39\,gの物質の体積を求めなさい。}{5\,cm$^3$}

\Instruction{\strut}
\QQ{金属をみがいたときに現れる、特有のかがやきを何というか。}{金属光沢}
\QQ{メスシリンダーでは、最小目盛りの何分の1まで目分量で読み取るか。}{10分の1}
\QQ{ガスバーナーで、炎に混じる空気の量を調節するねじを何というか。}{空気調節ねじ}
\QQ{金属をたたいて薄く広げられる性質を何というか。}{展性}
\QQ{上皿てんびんの分銅をのせたり取ったりするときに使う器具は何か。}{ピンセット}
